\documentclass[10pt,a4paper]{amsart}
\usepackage{amsmath,amssymb,amsthm,graphicx}
% Retained preamble for the Stacks Project chapter snapshot in this directory.
% The chapter's own preamble.tex is a copy of the shared Stacks Project preamble; the block
% below is sliced MECHANICALLY, line for line, from it. Nothing was added, removed or
% redefined; the wrapper class, the xr-hyper/hyperref wiring and the externaldocument list
% are omitted because the assembler supplies the document class.
%
% lines taken from preamble.tex: 3, 5, 9, 11, 12, 156, 157, 158, ... (51 lines total)
%
% package coverage restored from the chapter's own preamble (2026-10-13): xy, xy-2cell

\usepackage[all]{xy}\usepackage{verbatim}
\xyoption{2cell}
\UseAllTwocells

\newenvironment{reference}{\comment}{\endcomment}
\newenvironment{slogan}{\comment}{\endcomment}
\newenvironment{history}{\comment}{\endcomment}
\theoremstyle{plain}
\newtheorem{theorem}[subsection]{Theorem}
\newtheorem{proposition}[subsection]{Proposition}
\newtheorem{lemma}[subsection]{Lemma}
\theoremstyle{definition}
\newtheorem{definition}[subsection]{Definition}
\newtheorem{example}[subsection]{Example}
\newtheorem{exercise}[subsection]{Exercise}
\newtheorem{situation}[subsection]{Situation}
\theoremstyle{remark}
\newtheorem{remark}[subsection]{Remark}
\newtheorem{remarks}[subsection]{Remarks}
\numberwithin{equation}{subsection}
\def\lim{\mathop{\mathrm{lim}}\nolimits}
\def\colim{\mathop{\mathrm{colim}}\nolimits}
\def\Spec{\mathop{\mathrm{Spec}}}
\def\Hom{\mathop{\mathrm{Hom}}\nolimits}
\def\Ext{\mathop{\mathrm{Ext}}\nolimits}
\def\SheafHom{\mathop{\mathcal{H}\!\mathit{om}}\nolimits}
\def\SheafExt{\mathop{\mathcal{E}\!\mathit{xt}}\nolimits}
\def\Sch{\mathit{Sch}}
\def\Mor{\mathop{\mathrm{Mor}}\nolimits}
\def\Ob{\mathop{\mathrm{Ob}}\nolimits}
\def\Sh{\mathop{\mathit{Sh}}\nolimits}
\def\NL{\mathop{N\!L}\nolimits}
\def\CH{\mathop{\mathrm{CH}}\nolimits}
\def\proetale{{pro\text{-}\acute{e}tale}}
\def\etale{{\acute{e}tale}}
\def\QCoh{\mathit{QCoh}}
\def\Ker{\mathop{\mathrm{Ker}}}
\def\Im{\mathop{\mathrm{Im}}}
\def\Coker{\mathop{\mathrm{Coker}}}
\def\Coim{\mathop{\mathrm{Coim}}}
\DeclareMathSymbol{\boxtimes}{\mathbin}{AMSa}{"02}
\def\QCohstack{\mathcal{QC}\!\mathit{oh}}
\def\Cohstack{\mathcal{C}\!\mathit{oh}}
\def\Spacesstack{\mathcal{S}\!\mathit{paces}}
\def\Quotfunctor{\mathrm{Quot}}
\def\Hilbfunctor{\mathrm{Hilb}}
\def\Curvesstack{\mathcal{C}\!\mathit{urves}}
\def\Polarizedstack{\mathcal{P}\!\mathit{olarized}}
\def\Complexesstack{\mathcal{C}\!\mathit{omplexes}}
\def\Pic{\mathop{\mathrm{Pic}}\nolimits}
\def\Picardstack{\mathcal{P}\!\mathit{ic}}
\def\Picardfunctor{\mathrm{Pic}}
\def\Deformationcategory{\mathcal{D}\!\mathit{ef}}
\usepackage{cleveref}
\crefname{originallemma}{original lemma}{original lemmas}
\Crefname{originallemma}{Original lemma}{Original lemmas}
\makeatletter
\newlabel{properties-proposition-characterize-ample}{{01Q3}{source}{}{sourcepointer-1}{}}
\newlabel{properties-proposition-characterize-ample@cref}{{[originallemma][0][]01Q3}{[1][0][]source}{}{}{}}
\makeatother
\crefname{originallemma}{original lemma}{original lemmas}
\Crefname{originallemma}{Original lemma}{Original lemmas}
\makeatletter
\newlabel{spaces-divisors-definition-relatively-ample}{{0D31}{source}{}{sourcepointer-2}{}}
\newlabel{spaces-divisors-definition-relatively-ample@cref}{{[originallemma][0][]0D31}{[1][0][]source}{}{}{}}
\makeatother
\crefname{originallemma}{original lemma}{original lemmas}
\Crefname{originallemma}{Original lemma}{Original lemmas}
\makeatletter
\newlabel{spaces-limits-lemma-directed-colimit-finite-type}{{0829}{source}{}{sourcepointer-3}{}}
\newlabel{spaces-limits-lemma-directed-colimit-finite-type@cref}{{[originallemma][0][]0829}{[1][0][]source}{}{}{}}
\makeatother
\title{Stacks, Derived Categories of Spaces: resolution-property-ample-relative}
\author{The Stacks Project Authors}
\date{stacks-project:spaces-perfect.tex:local-snapshot-20261012; GNU FDL 1.2 or later}
\begin{document}
\maketitle
\noindent\textit{Editorial note.} Statement span is lines 7129-7140 of sources/stacks-spaces-perfect/pinned-original-source.tex; proof span is lines 7142-7174, the complete author proof (31 body lines).\par\medskip
\noindent\textit{Editorial note.} Navigation stubs: 3 cited-but-not-included labels are supplied as non-load-bearing pointers with their pinned tags where the snapshot tag index has them.\par\medskip
\noindent\textit{Editorial note.} extra\_packages is empty on purpose: the retained preamble already loads xy with options and a second declaration clashes in the assembler wrapper.\par\medskip
\noindent\textit{Editorial note.} Text is retained verbatim, including the author's own typos and placeholders.\par\medskip
\noindent\textit{Editorial note.} Queue index 50 of queue\_spaceperfect.json (pb=1368, stmt 7129-7140, proof 7142-7174).\par\medskip
\noindent\textit{Editorial note.} stacks-spaces-perfect wave sp01 (2026-10-14): fresh adjudication accept per lane-s3/sp01\_judgments; pool per Chief 12:4x unfreeze (reserve-pool workflow, six-check clean); spaces-perfect.tex 126 with-proof rows, 70 b107-occupied labels avoided, 56 free rows judged 32 accept / 24 reject.\par\medskip
\noindent\textit{Editorial note.} Cross-chapter same-label warning (NOT an occupancy conflict, Chief ruling 2026-10-13 evening): label 'lemma-resolution-property-ample-relative' also exists in stacks-perfect — a different proposition in a different chapter file.\par\medskip
\noindent\textit{Allowed standard prerequisites (editorial).}\begin{itemize}
\item Explicit prerequisite boundary. The cited inputs are prerequisites; the argument in the retained span is not replaced by any cited result.
\end{itemize}\medskip
\ifdefined\hypertarget\hypertarget{sourcepointer-1}{}\fi
\noindent\textit{Source reference.} Original lemma 01Q3: sources/stacks-spaces-perfect/pinned-original-source.tex, source lines 7174--7174. Cited by the retained statement or proof but not part of the collected spans; supplied only as a navigation pointer. The argument does not depend on its content.\par\medskip
\ifdefined\hypertarget\hypertarget{sourcepointer-2}{}\fi
\noindent\textit{Source reference.} Original lemma 0D31: sources/stacks-spaces-perfect/pinned-original-source.tex, source lines 7174--7174. Cited by the retained statement or proof but not part of the collected spans; supplied only as a navigation pointer. The argument does not depend on its content.\par\medskip
\ifdefined\hypertarget\hypertarget{sourcepointer-3}{}\fi
\noindent\textit{Source reference.} Original lemma 0829: sources/stacks-spaces-perfect/pinned-original-source.tex, source lines 7174--7174. Cited by the retained statement or proof but not part of the collected spans; supplied only as a navigation pointer. The argument does not depend on its content.\par\medskip
\begin{lemma}
\label{lemma-resolution-property-ample-relative}
Let $S$ be a scheme.
Let $f : X \to Y$ be a morphism of algebraic spaces over $S$. Assume
\begin{enumerate}
\item $Y$ is quasi-compact and quasi-separated and has the resolution property,
\item there exists an $f$-ample invertible module on $X$
(Divisors on Spaces, Definition
\ref{spaces-divisors-definition-relatively-ample}).
\end{enumerate}
Then $X$ has the resolution property.
\end{lemma}

\begin{proof}
Let $\mathcal{F}$ be a finite type quasi-coherent $\mathcal{O}_X$-module.
Let $\mathcal{L}$ be an $f$-ample invertible module.
Choose an affine scheme $V$ and a surjective \'etale morphism $V \to Y$.
Set $U = V \times_Y X$. Then $\mathcal{L}|_U$ is ample on $U$.
By Properties, Proposition
\ref{properties-proposition-characterize-ample}
we know there exists finitely many maps
$s_i : \mathcal{L}^{\otimes n_i}|_U \to \mathcal{F}|_U$
which are jointly surjective. Consider the quasi-coherent
$\mathcal{O}_Y$-modules
$$
\mathcal{H}_n =
f_*(\mathcal{F} \otimes_{\mathcal{O}_X} \mathcal{L}^{\otimes n})
$$
We may think of $s_i$ as a section over $V$ of the sheaf
$\mathcal{H}_{-n_i}$. Suppose we can find finite locally
free $\mathcal{O}_Y$-modules $\mathcal{E}_i$ and maps
$\mathcal{E}_i \to \mathcal{H}_{-n_i}$ such that
$s_i$ is in the image. Then the corresponding maps
$$
f^*\mathcal{E}_i
\otimes_{\mathcal{O}_X}
\mathcal{L}^{\otimes n_i} \longrightarrow \mathcal{F}
$$
are going to be jointly surjective and the lemma is proved. By
Limits of Spaces, Lemma \ref{spaces-limits-lemma-directed-colimit-finite-type}
for each $i$ we can find a finite type quasi-coherent submodule
$\mathcal{H}'_i \subset  \mathcal{H}_{-n_i}$
which contains the section $s_i$ over $V$.
Thus the resolution property of $Y$ produces surjections
$\mathcal{E}_i \to \mathcal{H}'_i$ and we conclude.
\end{proof}

\end{document}
